La crescente disponibilità di dati personali e la diffusione di questi su una grande quantità di piattaforme pubbliche rende tecniche tradizionali di anonimizzazione inadeguate a garantire la protezione della privacy.
La privacy differenziale (DP) costituisce un framework matematico che, tramite l'utilizzo di meccanismi che agiscono aggiungendo quantità calibrate di rumore alle operazioni effettuate sui dati, offre garanzie formali sulla riservatezza dei dati anche in presenza di attaccanti dotati di informazioni ausiliarie.

In questo lavoro di tesi si discutono i fondamenti teorici della DP, analizzando alcuni meccanismi differenzialmente privati, il significato operativo dei parametri epsilon e delta e le proprietà del budget di privacy. A partire da questo studio si sviluppa GoDP, un'applicazione scritta in Go basata su Privacy on Beam pensata per consentire la creazione di distribuzioni di dati DP a partire da dataset arbitrari. Lo strumento è ideato per abbassare la barriera di ingresso alla DP, permettendo a utenti non esperti di applicare la DP a distribuzioni di dati sensibili.

Viene valutata la validità dello strumento tramite un esperimento basato sulla simulazione di un differencing attack, ipotizzando un attaccante con determinate caratteristiche e applicando il modello di DP centralizzato. I risultati mostrano come, senza le garanzie offerte dalla DP, l'attacco riesca ad identificare il dato sensibile, mentre con GoDP viene aggiunto rumore calibrato rendendo l'inferenza impossibile o imprecisa a seconda dei parametri del meccanismo, evidenziando il compromesso tra garanzie di privacy e utilità dei dati proprio della DP.